\documentclass[11pt,a4paper]{article}

\usepackage[utf8]{inputenc}
\usepackage[T1]{fontenc}
\usepackage{lmodern}
\usepackage{geometry}
\usepackage{hyperref}
\usepackage{setspace}
\usepackage{enumitem}
\usepackage{parskip}
\usepackage{graphicx}
\usepackage{float}
\usepackage{subcaption}

\geometry{margin=2.5cm}
\setstretch{1.08}
\hypersetup{hidelinks}
\graphicspath{{../results/}{results/}}

\title{Rapport interm\'ediaire~-- \'Etat d'avancement}
\author{}
\date{6 avril 2026}

\begin{document}
\maketitle

\section{Ce qui fonctionne}

\subsection{Ordonnancement coop\'eratif FIFO}

L'ordonnanceur est op\'erationnel utilisant actuellement la politique FIFO. La \texttt{ready\_queue} est une
\texttt{TAILQ} BSD~: insertion en $O(1)$ en queue, extraction en $O(1)$ en t\^ete.

\texttt{thread\_yield} place le thread courant en \'etat \texttt{THREAD\_READY}
en queue puis bascule vers le premier thread disponible via \texttt{swapcontext}.
\texttt{thread\_yield\_to} cible un thread pr\'ecis~; si celui-ci n'est pas pr\^et,
il se rabat sur \texttt{thread\_yield}.

\texttt{thread\_join} impl\'emente une d\'etection de deadlock par parcours de la
cha\^ine \texttt{joined\_by} avant de bloquer (complexité linéaire en la taille de la chaine, à améliorer si possible). Si la cible est d\'ej\`a termin\'ee,
le join est imm\'ediat, sans changement de contexte. Sinon le thread courant passe
en \texttt{THREAD\_BLOCKED} et c\`ede la main.

\texttt{thread\_exit} marque le thread \texttt{THREAD\_TERMINATED}, d\'epose la
valeur de retour, puis r\'eveille le joineur (\texttt{joined\_by}) s'il en existe
un, prend le premier thread pr\^et disponible sinon, ou bascule sur la
\textit{cleanup stack} si aucun thread n'est disponible.

Tous les tests fournis (01 \`a 81) passent. Valgrind ne signale aucune fuite
m\'emoire (avec les options \texttt{-s --leak-check=full --show-reachable=yes --track-origins=yes}).

\subsection{Gestion des piles~: pool et guard pages}

Chaque pile est allou\'ee via \texttt{mmap} avec une \textit{guard page} de 4~Ko
en bas de pile (\texttt{mprotect(PROT\_NONE)}), ce qui d\'etecte les d\'ebordements
d\`es le premier acc\`es hors-limites. Chaque pile est enregistr\'ee aupr\`es de
Valgrind via \texttt{VALGRIND\_STACK\_REGISTER}.

Plut\^ot que de lib\'erer syst\'ematiquement les piles \`a chaque terminaison de
thread, un \textbf{pool de piles} conserve jusqu'\`a \texttt{MAX\_POOL\_STACKS}. entr\'ees r\'eutilisables.
\texttt{stack\_pool\_alloc} d\'epile une entr\'ee existante en $O(1)$ au lieu de
rappeler \texttt{mmap}. Ce m\'ecanisme r\'eduit significativement le co\^ut de
cr\'eation lors des tests \`a grande \'echelle (\texttt{22-create-many},
\texttt{23-create-many-recursive})~: la diff\'erence est mesurable d\`es quelques
centaines de threads.

\subsection{Garbage collection des zombies}

Un thread termin\'e non encore join\'e est plac\'e dans une \texttt{zombie\_queue}~:
sa m\'emoire n'est pas lib\'er\'ee imm\'ediatement car un futur \texttt{thread\_join}
doit encore lire \texttt{retval}. \`A la fin du programme, \texttt{do\_final\_cleanup} lib\`ere tous les zombies r\'esiduels et vide le
pool. Cette fonction s'ex\'ecute sur une pile neutre de 8~Ko (\texttt{cleanup\_stack})
pour ne jamais lib\'erer la pile sur laquelle on s'ex\'ecute.

\subsection{Mutex}

Les mutex disposent d'une \texttt{waiting\_queue} par instance.
\begin{itemize}[noitemsep]
  \item \textbf{Fast path ($O(1)$)~:} si le mutex est libre,
        \texttt{thread\_mutex\_lock} l'acquiert sans changement de contexte.
  \item \textbf{Slow path~:} si le mutex est occup\'e, le thread courant passe en
        \texttt{THREAD\_BLOCKED} et est ins\'er\'e dans la \texttt{waiting\_queue}.
        \texttt{thread\_mutex\_unlock} transfère la propri\'et\'e directement au
        premier waiter (\texttt{locked} reste \`a 1). Le thread r\'eveill\'e est remis en \texttt{ready\_queue}
        en \'etat \texttt{THREAD\_READY}.
\end{itemize}

Si la pr\'eemption est active (\texttt{PREEM\_ENABLED}), les sections critiques
du mutex sont prot\'eg\'ees par \texttt{preem\_block}/\texttt{preem\_unblock}.

\subsection{Pr\'eemption}

La pr\'eemption repose sur \texttt{SIGVTALRM} et
\texttt{setitimer(ITIMER\_VIRTUAL)}. L\`alarme est d\'eclench\'ee toutes les 5~ms (configurable via \texttt{init\_prem}).
Le gestionnaire de signal (\texttt{preemption\_handler}) appelle
\texttt{thread\_yield}, ins\'erant le thread courant en fin de queue et basculant
vers le prochain thread pr\^et.

Toutes les sections critiques du planificateur (cr\'eation, yield, join, mutex)
sont prot\'eg\'ees par \texttt{preem\_block}/\texttt{preem\_unblock}
(\texttt{sigprocmask(SIG\_BLOCK\-/SIG\_UNBLOCK)}) pour \'eviter toute r\'eentrance
asynchrone. L'option \texttt{SA\_RESTART} est pass\'ee \`a \texttt{sigaction} pour
que les appels syst\`eme interrompus soient relanc\'es automatiquement.

La pr\'eemption est activ\'ee \`a la compilation avec \texttt{-DENABLE\_PREEMPTION}
pour ne pas p\'enaliser les tests qui n'en ont pas besoin. Le test
\texttt{71-preemption} passe~: un thread est bien interrompu et un 
autre progress.

\section{Performances~: comparaison avec pthreads}

Les premiers r\'esultats montrent un avantage net de notre biblioth\`eque sur
les tests orient\'es cr\'eation massive de threads et synchronisation coop\'erative.
Les figures ci-dessous illustrent les mesures obtenues avec le script de benchmark.

\begin{figure}[H]
      \centering
      \begin{subfigure}[t]{0.48\linewidth}
            \centering
            \includegraphics[width=\linewidth]{21-create-many.png}
            \caption{\texttt{21-create-many}}
            \label{fig:create-many}
      \end{subfigure}
      \hfill
      \begin{subfigure}[t]{0.48\linewidth}
            \centering
            \includegraphics[width=\linewidth]{22-create-many-recursive.png}
            \caption{\texttt{22-create-many-recursive}}
            \label{fig:create-many-rec}
      \end{subfigure}

      \vspace{0.6em}

      \begin{subfigure}[t]{0.48\linewidth}
            \centering
            \includegraphics[width=\linewidth]{23-create-many-once.png}
            \caption{\texttt{23-create-many-once}}
            \label{fig:create-many-once}
      \end{subfigure}
      \hfill
      \begin{subfigure}[t]{0.48\linewidth}
            \centering
            \includegraphics[width=\linewidth]{61-mutex.png}
            \caption{\texttt{61-mutex}}
            \label{fig:mutex}
      \end{subfigure}

      \caption{Comparaison des temps d'ex\'ecution entre notre impl\'ementation et \texttt{pthread}.}
      \label{fig:perf-grid}
\end{figure}

Ces r\'esultats s'expliquent par la nature des tests~:
\begin{itemize}[noitemsep]
  \item \textbf{Cr\'eation massive de threads (21/22/23).} Notre biblioth\`eque
        cr\'ee des threads utilisateur plus l\'egers, avec r\'eutilisation de
        piles via le pool. \`A l'inverse, \texttt{pthread} implique davantage de
        co\^uts noyau et de synchronisations internes pour chaque cr\'eation.
  \item \textbf{Commutation de contexte en espace utilisateur.} Les
        basculements se font principalement via \texttt{swapcontext}, sans
        ordonnanceur noyau \`a chaque changement de thread, ce qui r\'eduit le
        co\^ut moyen par commutation.
  \item \textbf{Mutex en mode coop\'eratif (61-mutex).} Le r\'eveil/transfert de
        propri\'et\'e est direct entre threads utilisateur (file d'attente locale),
        avec moins de transitions syst\`eme dans ce sc\'enario monocoeur.
\end{itemize}

En pratique, sur ces charges orient\'ees ``beaucoup de threads + beaucoup de
commutations'', notre impl\'ementation est syst\'ematiquement devant \texttt{pthread}.
Ce n'est pas une sup\'eriorit\'e universelle~: sur des charges fortement
calculatoires exploitant le parall\'elisme mat\'eriel, \texttt{pthread} peut
redevenir meilleur tant que notre version multicoeur n'est pas activ\'ee.


\section{Limites connues}

\begin{itemize}[noitemsep]
  \item \textbf{Boucle semi-active dans \texttt{thread\_join}~:} l'attente utilise
        une boucle \texttt{while} avec \texttt{swap\_thread}. Dans des cas avec de
        nombreux threads bloqu\'es simultan\'ement, cela peut g\'en\'erer des
        changements de contexte inutiles. Une attente purement passive (r\'eveil
        uniquement d\'eclench\'e par \texttt{thread\_exit}) serait plus propre~;
        c'est pr\'evu pour la suite.
  \item \textbf{Graphes de performance incomplets~:} le script \texttt{make graphs}
        est fonctionnel mais les courbes doivent encore \^etre produites pour
        l'ensemble des tests $\geq 21$ avec variation du nombre de threads.
\end{itemize}

\section{Suite du projet}

\subsection{Support multicoeur}

Utiliser plusieurs threads noyau internes (pthreads) pour ex\'ecuter les threads
utilisateur en parall\`ele. La \texttt{ready\_queue} et les structures de
l'ordonnanceur devront \^etre prot\'eg\'ees par des verrous
(\texttt{pthread\_spinlock\_t}). On mesurera si le gain de parall\'elisme compense
le surco\^ut de synchronisation, et pour quel type de workload.

\subsection{Signaux inter-threads}

Permettre l'envoi de signaux entre threads utilisateur ind\'ependamment des signaux
syst\`eme (d\'ej\`a utilis\'es pour la pr\'eemption), et potentiellement impl\'ementer plusieurs fonctionnalités
\texttt{sigwait} pour qu'un thread dorme jusqu'\`a la r\'eception d'un signal.

\subsection{Am\'elioration du \texttt{thread\_join}}

Passer \`a une attente purement passive~: le thread joineur est d\'epos\'e dans
\texttt{joined\_by} et ne reprend la main que lorsque \texttt{thread\_exit} le
r\'eveille explicitement, sans boucle de polling.

\subsection{Finalisation des graphes de performance}

Compléter le script \texttt{make graphs} pour couvrir tous les tests $\geq 21$
avec variation du nombre de threads, et produire des courbes propres
(comparaison biblioth\`eque\,/\,pthreads) pour le rapport final.

\end{document}